% Data and Method v2 (20 Nov 2026). Values checked against configs/*.yaml and scripts (splits, training schedule, XGBoost settings, demeaning).
% TODO fill: Argo GDAC access date and DOI for the earlier months. TODO: all \cite keys must be verified by the author with DOIs.

\section{Data and method}
\label{sec:data}
\emph{Data.} All fields are on a daily $0.25^\circ$ grid over $45$--$105^\circ$E and $5$--$30^\circ$N ($101\times241$ cells). The target is GLORYS12 potential temperature \cite{glorys12} at 15 depths from 0 to 1000~m (0, 5, 10, 20, 30, 50, 75, 100, 125, 150, 200, 300, 500, 700 and 1000~m). The inputs are OSTIA reprocessed SST (v202003) \cite{ostia} and DUACS multi-mission sea level anomaly (SLA, v202411) \cite{duacs}. Of the 12,038 ocean cells, 9,140 are valid at all 15 depths (the \emph{deep set}); every headline number uses this set. From July 2021 GLORYS12 is the interim production of the same product (\texttt{cmems\_mod\_glo\_phy\_my\_0.083deg\_P1D-m}, v202311), so Tests~B and C use an interim target.

\emph{Protocol.} Splits are chronological: training 2010-02-04 to 2017-12-31; validation 2018-02-01 to 2020-01-31 (widened from one year on 25 October 2026, before any model was trained); Test~A 2020-03-02 to 2021-06-27; Test~B 2021-08-01 to 2025-12-31; and Test~C 2026-01-01 to 2026-05-15 (135 days). Gaps between training, validation, Test~A and Test~B are at least 30 days; Test~C follows Test~B directly. Test~C is prospective: its protocol and four predictions were committed on 15 November 2026 (commit 9869230) before any 2026 file was downloaded, although the 2026 data existed publicly before the models were frozen. All models were frozen before any test period was opened. Each test period was scored once under a lock file; the only reruns were for code or output fixes and were logged, including an Argo scoring of Test~C that returned no rows because the profile-pair file lacked 2026 profiles.

\emph{Argo.} Delayed-mode and adjusted Argo profiles are matched to the nearest $0.25^\circ$ cell and day, converted to potential temperature with TEOS-10 \cite{teos10}, and compared at the same cell, day and depth with every estimate and with GLORYS12; 0~m is dropped. The 2026 profiles were downloaded on 16 November 2026 through argopy \textbf{[FILL: access date and DOI for the earlier months]}; Test~C has 507 such profiles. Argo is assimilated by GLORYS12, so it is not independent of the target.

\emph{Models.} Inputs and targets are anomalies relative to a per-cell harmonic seasonal cycle (three harmonics) fitted on training days only; inputs are divided by training standard deviations, and the network predicts the target anomaly (scaled by 2). The U-Net \cite{unet} has width 32 and about 0.47 million parameters. It maps SST and SLA anomalies and an ocean mask to 15 anomaly maps and has no season input. Training used AdamW (learning rate $10^{-3}$, weight decay $10^{-4}$), cosine annealing over at most 25 epochs, batch size 8, every second training day (1,444 days) and early stopping on the validation loss (patience 6). The \emph{frozen ensemble} averages six U-Nets: three seeds with same-day inputs and three with inputs at lags of 0, 7, 14 and 30 days. The \emph{demeaned ensemble} is identical except that the daily deep-set mean of SLA is removed before standardisation; it was pre-registered as co-primary and is always reported.

\emph{Baselines.} (i)~Seasonal climatology. (ii)~A per-cell, per-depth linear regression on the SST and SLA anomalies. (iii)~XGBoost \cite{xgboost}, one model per depth, with the two anomalies, latitude, longitude and the sine and cosine of the day of year as features and fixed untuned settings (300 trees, learning rate 0.05, depth 6), fitted on every tenth training day and 1,500 random deep-set cells per day; it covers the deep set only. (iv)~A single U-Net (seed 0), which is the first member of the frozen ensemble. All are fitted on training days only.

\emph{Argo-calibrated correction.} For each of the 14 Argo depth levels, the mean of (estimate minus Argo) over validation-window profiles is subtracted from the estimate at every cell and day: one constant per depth, fitted on validation Argo only. The rule was pre-registered (commit 035f117) after the test periods had been scored, so it is pre-specified but not blind. Its extension to all estimates (commit 3d1ed25) and the comparison of XGBoost and the ensemble by season are exploratory.

\emph{Metrics.} Against GLORYS12, pooled RMSE is the square root of the mean over the 15 depths of the per-depth MSE on the deep set. Against Argo it is the RMSE over all matched profile-depth values at 14 levels, on the rows where every estimate is finite. Intervals for paired RMSE differences use a moving-block bootstrap over days (30-day blocks, 2,000 resamples); they reflect day-sampling variability, not training randomness, and Test~C contains about 4.5 blocks.
